\section{Methods}
\label{sec:methods}

\subsection{Design, reporting and protocol status}
This is a retrospective, pre-deployment evaluation of fixed models at two
sites. A machine-readable protocol (registry
\texttt{protocols/C3E6\_stage2}, version 0.1.0) was frozen on 18 July 2026,
before any MIMIC data were accessed. Five dated amendments followed (versions
0.2.0--0.4.2, 28 July to 4 August 2026). All of them came before any external
result existed, and the last came after the image-only model had been trained
and its internal validation AUROC seen. External point estimates were first
inspected on 7 August 2026 (08:15, UTC+6). Several analysis decisions were
taken after that moment and are identified as post hoc in
\cref{tab:status} and in the protocol history at the end of the paper. The
audit and every analysis labelled \emph{Revision R1} were carried out on 6
October 2026. We follow TRIPOD+AI reporting where it applies to an evaluation
of existing models.

\subsection{Sites, data and cohorts}
The source site was MIMIC-CXR-JPG v2.1.0~\cite{johnson2019mimiccxrjpg,johnson2019mimiccxr}
(Beth Israel Deaconess Medical Center). The external site was CheXpert
Plus~\cite{chambon2024chexpertplus} (Stanford). It was used only for evaluation;
the registry prohibits its use for architecture, early-stopping, temperature,
threshold, confidence-score, abstention-cutoff, text-cleaning or context-rule
selection. \Cref{tab:cohort} gives the source exclusion waterfall, the
patient-disjoint tiers and the external cohort, all taken from the committed
cohort artifacts.

Report sections were extracted with a curated header vocabulary rather than by
position. 66 reports had no recognised header and were excluded. 134 reports
merged findings and impression into one section; 133 further studies were
removed for this reason in the sequential waterfall, so 199 studies were removed
in total, not the 200 stated in the previous draft. The prespecified
evaluation tier (5,000 patients) and the calibration tier (3,000 patients) were
assigned from the official training split using patient identifiers and a
fixed seed alone. All tiers are patient-disjoint, and the official validate and
test splits were not analysed. At the external site, one study has two frontal
images whose permitted-context fields differ: one image qualifies as
informative and one does not. The study is therefore counted in two natural
states, and only its informative image enters the intervention cohort.
One external image matches its published checksum but does not decode and was
excluded.

\begin{table}[t]
\centering\small
\caption{Cohort construction. Source rows are sequential exclusions from the
committed Stage 3C waterfall; tiers are from the Stage 3E partition; external
counts are from the Stage 9b artifact.}
\label{tab:cohort}
\begin{adjustbox}{max width=\linewidth}
\begin{tabular}{lrr}
\toprule
Step (source site) & Removed & Remaining\\
\midrule
All reports/studies & -- & 227,835\\
No recognised section header & 66 & 227,769\\
Findings and impression merged & 133 & 227,636\\
No frontal (AP/PA) image & 9,688 & 217,948\\
Official split assigned & 0 & 217,948\\
Context not informative ($<3$ effective tokens) & 15,001 & 202,947\\
No impression section & 29,375 & 173,572\\
\midrule
Tier & Patients / studies / frontal images & \\
Model training & 49,959 / 146,253 / 162,823 & \\
Threshold calibration & 3,000 / 9,233 / 10,264 & \\
Prespecified evaluation (primary) & 5,000 / 14,766 / 16,456 & \\
Official validate (not analysed) & 448 / 1,391 / 1,584 & \\
Official test (not analysed) & 282 / 1,929 / 2,155 & \\
\midrule
External site (CheXpert Plus) & & \\
Usable frontal images / studies / patients & 191,070 / 187,673 / 64,709 & \\
Intervention cohort (N1 at $T=3$): images / studies / patients & 135,561 / 132,923 / 53,110 & \\
\bottomrule
\end{tabular}

\end{adjustbox}
\end{table}

\subsection{Pre-diagnostic context}
Permitted context at the source was the INDICATION and HISTORY sections, with
registered synonyms. At the external site it was the published
\texttt{section\_clinical\_history} and \texttt{section\_history} fields,
concatenated. 17,551 source reports (7.70\%) contain a WET READ (preliminary
interpretation) section, and 181 place a pre-diagnostic section after a
diagnostic one. Header-driven extraction keeps both out of the context field;
post-diagnostic sections (WET READ, provisional findings, recommendation,
notification, addendum, comment) are never extracted. Section headers do not
prove that the text existed before the image was read. We therefore inspected
extracted context. No extracted context at either site contained a
report-section header (\ScanHeaderSrc{} and \ScanHeaderExt{}). Conclusion-style phrases
(e.g.\ ``consistent with'', ``unchanged from'') occurred in \ScanConclSrc{} and
\ScanConclExt{}, and comparison language in \ScanComparSrc{} and
\ScanComparExt{}. Context named one of the target findings or a close synonym
in \ScanTargetSrc{} of source and \ScanTargetExt{} of external studies, often as
a suspected or known condition (``evaluate progression of right effusion'').
That is genuinely pre-diagnostic for the current study but label-informative,
which is why M2 is reported as a shortcut control. Report text was inspected only in the protected environment and
appears in no artifact.

\subsection{Context states and interventions}
An \emph{effective token} is a whitespace-delimited token containing an
alphanumeric character after de-identification placeholder runs are removed. A
study is informative (N1) if its context has at least $T$ effective tokens
($T=3$ primary, registered sensitivities $T=5,10$), low-information (N2) if it
has between 1 and $T-1$, and absent (N3) if it has none. The rule was fixed on
the source and applied unchanged externally. Interventions apply to N1
studies: C0 leaves context unchanged; C1 replaces it with the token
\texttt{[NO\_CONTEXT]}; C2 replaces it with another patient's context. The C2
replacement is a deterministic rotation within strata defined by split and
context-length bin (bins 0--4, 5--9, 10--19, 20--39, 40--79, $\ge$80 effective
tokens), seed 20260718. Within each stratum it walks forward past donors from
the same patient.

Two implementation facts were established by the audit
(\cref{sec:res-c2}) and are stated here so the interventions are described as
run. First, C2 was applied to \emph{image} rows, not studies, so images of one
study can receive context from different donor studies. This happened for
\CTwoMultiDonorSrc{} of \CTwoMultiImgSrc{} multi-image source studies and
\CTwoMultiDonorExt{} of \CTwoMultiImgExt{} external ones. Second, the C2 donor
realisation depends on the cohort being permuted. Changing the cohort, as the
label-source and $T$ sensitivities do, changes every recipient's donor.

\subsection{Labels}
The targets were cardiomegaly, edema, consolidation, atelectasis and pleural
effusion. Source labels were generated by a port of CheXbert~\cite{smit2020chexbert}
(checkpoint SHA-256 \texttt{6550703c\ldots0fde1} from
\texttt{StanfordAIMI/RRG\_scorers}; upstream commit \texttt{6d22a96})
applied separately to the impression and findings sections. The port's fidelity
check covered 4 public sample reports (56 cells, all matching). External labels
are the CheXbert labels distributed with CheXpert Plus; the checkpoint and
pre-processing used to produce them are not documented in the release, and we
did not regenerate them. The registry fixes impression-derived labels as
primary and findings-derived labels as a mandatory sensitivity (S1), and
forbids full-report labels because the full report contains the context
fields.

Values are positive (1), negative (0), uncertain ($-1$) or unmentioned (blank).
Only positive and negative cells carry supervision or contribute to error.
Uncertain and unmentioned cells are masked; neither is mapped to negative.
Source labels are study-level. External labels are image-level, and a study
takes the first observed (0/1) value across its frontal images. A source study
enters the impression cohort only if its report has an impression section, and
the findings cohort only if it has a findings section. External studies are
retained whatever sections they have, with missing labels masked. This asymmetry
is central to \cref{sec:res-label-source}.

\subsection{Models and training}
M1 is DenseNet-121 (ImageNet-initialised) with five sigmoid outputs. M2 is
BERT-base-uncased (128 tokens, CLS pooling, fine-tuned) with five sigmoid
outputs. M3 averages the M1 and M2 probabilities and has no trainable
parameters. M4 concatenates the pooled DenseNet-121 features (1,024) with the BERT
CLS embedding (768), followed by Linear(1792, 512), ReLU, dropout 0.2 and
Linear(512, 5), trained end to end. Images were converted to greyscale, resized
once to $224\times224$ by bilinear interpolation without preserving aspect
ratio, stored as JPEG (quality 95), replicated to three channels and normalised
with ImageNet statistics. The pipeline was identical at both sites, and no
augmentation was used.

Training used the model-train tier only (146,253 studies, 49,959 patients). A
patient-disjoint 5\% internal-validation split (2,498 patients, 7,919 images,
fixed seed) was used for early stopping. The image is the training unit. The
loss is masked binary cross-entropy,
\begin{equation}
\mathcal{L}=\frac{\sum_{b\in\mathcal{B}}\sum_{j=1}^{5} m_{bj}\,
\ell\!\left(z_{bj},y_{bj}\right)}{\max\!\left(1,\sum_{b\in\mathcal{B}}\sum_{j} m_{bj}\right)},
\end{equation}
over the images $b$ of a batch, with logit $z$, label $y$ and mask $m$. It is a
micro-average over labelled cells with no class or study weights. Images whose
five cells are all unlabelled contribute nothing. The optimiser was AdamW
(learning rate $10^{-4}$ for image and fusion parameters and $2\times10^{-5}$
for the text encoder, weight decay $10^{-4}$), with batch size 32 in fp32, at
most 15 epochs, and early stopping on internal-validation macro AUROC with
patience 3. The selected checkpoints were M1 epoch 3 (AUROC 0.877), M2 epoch 3
(0.763) and M4 epoch 4 (0.888). Under the second training seed (20260719), with
identical settings, the selected checkpoints were M1 epoch 3 (0.874), M4 epoch 1
(0.887) and M2 epoch 2 (0.768); the replicate M3 averages the replicate M1 and
M2. The committed Stage 6b audit artifact misreports
the M2 training curve by splicing epochs from a discarded run; the checkpoint
and its selection are unaffected (supplement).

\subsection{The frozen selective policy}
Study probabilities $p_{ij}$ are means over a study's frontal images. For
pathology $j$, the decision threshold $t_j$ maximises balanced accuracy over a
grid of 999 values, $k/1000$ for $k=1,\ldots,999$, on supervised
calibration-tier cells; ties go to the lowest threshold. The decision is
$\hat y_{ij}=\mathbb{1}[p_{ij}\ge t_j]$. Study confidence is
$g_i=\min_j|2p_{ij}-1|$, and the abstention cutoff $\tau$ is the 20th percentile
of $g$ over calibration studies (linear interpolation). A study is accepted when
$g_i\ge\tau$, so ties at the cutoff are accepted and coverage can exceed the
target, as with M2 under C1. All thresholds and cutoffs were selected on the
calibration tier (9,233 studies) and applied unchanged to the source evaluation
tier and the external site. Re-running calibration-tier inference reproduced
every frozen threshold and cutoff exactly (supplement). The confidence $g$ is
centred at $p=0.5$, whereas decisions are made at $t_j$, which ranged from
0.41 to 0.99 across models and pathologies. Whether $g$ ranks the correctness of the actual decisions is
therefore an empirical question (\cref{sec:res-selective}).

\subsection{Endpoint}
\label{sec:endpoint}
For study $i$ let $n_i=\sum_j m_{ij}$ be its number of labelled cells,
$w_i=\sum_j m_{ij}\mathbb{1}[\hat y_{ij}\ne y_{ij}]$ its number of wrong
labelled cells, and $a_i=\mathbb{1}[g_i\ge\tau]$ its acceptance. The registered
primary metric is $\mathrm{MET001}=\operatorname{mean}(L_i\mid g_i\ge\tau)$, with
$L_i$ the study's five-label Hamming error over supervised cells. The
original implementation computed
\begin{equation}
R_{\text{orig}}=\frac{\sum_i a_i\,w_i/\max(n_i,1)}{\sum_i a_i}.
\end{equation}
A study with $n_i=0$ therefore contributes error 0 to the numerator and 1 to
the denominator, so it is counted as correctly classified. The code's own
documentation states the opposite intent (``unsupervised cells are excluded
rather than counted correct''). We treat this as a verified implementation
defect and correct it by taking the registered mean over studies whose $L_i$
is defined:
\begin{equation}
R_{\text{eval}}=\frac{\sum_i a_i\,\mathbb{1}[n_i>0]\,w_i/n_i}{\sum_i a_i\,\mathbb{1}[n_i>0]},
\qquad R_{\text{orig}} = R_{\text{eval}}\cdot\Pr(n_i>0\mid a_i=1).
\end{equation}
The identity shows that the original estimator multiplies the error among
evaluable accepted studies by the fraction of accepted studies that are
evaluable. A cross-site difference in that fraction therefore appears as a
difference in risk. $R_{\text{eval}}$ is reported as the corrected primary and
$R_{\text{orig}}$ is retained. As sensitivity analyses we also report
cell-pooled risk, $\sum_i a_iw_i/\sum_i a_in_i$; pathology-macro risk, the
unweighted mean of the five per-pathology error rates over labelled accepted
cells; and pathology-specific risks, which the original SAP made the primary
reporting unit. Coverage is reported over all eligible studies, which is the
quantity a deployed policy controls, and over evaluable studies.
$R_{\text{eval}}$ is an error rate over cells that a labeller chose to label in
accepted studies. It is not a fully observed five-disease clinical error rate,
and evaluable studies are a selected subset whose composition differs between
sites.

\subsection{Hypotheses and decision rules}
Write $R^{s}_{m}(c)$ for selective risk at site $s\in\{S,E\}$ for model $m$
under condition $c$, at the frozen 80\% cutoff. The registered hypotheses
(\texttt{hypothesis\_registry.yaml}, 18 July 2026) are:
\begin{align*}
\text{H1 (confirmatory, } m\in\{\text{M3,M4}\}\text{):}\quad
&\Delta_m = R^{E}_{m}(\text{C0})-R^{S}_{m}(\text{C0}) > 0,\\
\text{H2 (confirmatory):}\quad
&\Omega_m = \bigl[R^{E}_{m}(\text{C1})-R^{E}_{m}(\text{C0})\bigr]-\bigl[R^{S}_{m}(\text{C1})-R^{S}_{m}(\text{C0})\bigr] > 0,\\
\text{H3 (confirmatory):}\quad
&\Omega_{m}-\Omega_{\text{M1}} > 0,\\
\text{H4 (secondary, each site):}\quad
&R_{m}(\text{C2})-R_{m}(\text{C1}) \ge 0 .
\end{align*}
H1 and H2 are confirmed when the 95\% interval lies above zero \emph{and} the
point estimate is at least 0.02 (registry \texttt{failure\_criteria}). The
registry gives H3 no materiality and gives H4 a non-negative direction with no
materiality. The pre-results implementation of H4 (commit \texttt{b6052a7}, 6
August 2026) declares it supported when the interval lies above zero. The
previous draft's statement that every hypothesis needed a 0.02 point estimate
was therefore wrong for H4. The H4 verdicts it reported, including
``confirmed'' for M4 at the source with \HFourMFoursourceOrigPt{}, were correct under H4's own
registered rule.

Two registered estimands were not implemented as written. M1 ignores text, so
$\Omega_{\text{M1}}\equiv 0$ and the registered H3 is algebraically identical to
H2 for the same model. The analysis code written on 7 August 2026, after the
external point estimates had been inspected, instead defined H3 as
$\Delta_m-\Delta_{\text{M1}}$ and applied the 0.02 materiality. That
operationalisation is reported here as post hoc. The same code also computed
an interaction under C2 that H2 does not register, and we report it as
descriptive. We report every contrast with one verdict function that
distinguishes five outcomes: confirmed; same direction but not material;
inconclusive (interval includes zero); opposite direction, not material; and
opposite direction, material. Separately, it records whether the interval
itself clears materiality and whether it lies inside $(-0.02,0.02)$. A point
estimate above 0.02 with an interval excluding zero does not show that the true
effect exceeds 0.02.

The 0.02 materiality threshold was registered without an external clinical
justification. We regard it as a convention, not a clinically validated
minimum important difference. Requiring the point estimate to reach it has an
elementary consequence. For an approximately normal estimate with standard
error $s$, the probability of confirmation when the true effect equals $\delta$
is $\Pr\{Z\ge\max(0,1.96-\delta/s)\}\le 1/2$, with equality once $\delta\ge
1.96s$. This is the familiar 50\% power at a decision boundary, not a new
methodological result. The study's power analysis was accordingly
based on the minimum detectable effect (0.021--0.027 at 80\% power on the
evaluation partition, under the registered grid of base risks, intra-patient
correlations and coverages), not on materiality.

\subsection{Statistical inference}
Intervals are 95\% percentile intervals from 2,000 patient-clustered bootstrap
replicates with seed 20260718. Within-site contrasts resample a site's patients
once per replicate and read every condition and model on that draw (paired).
Cross-site contrasts resample the two sites independently. H3 reads both models
from the same within-site draw. The revision reimplemented the bootstrap as
patient-weighted sums and replays the original random draws exactly: every
point estimate and interval in the original Stage 10 artifacts (primary, S1,
$T=5$, $T=10$, seed replicate; 63 rows) was reproduced exactly at the six-decimal precision of the artifacts.

The intervals condition on the fitted models, on the thresholds and cutoffs
selected from the one realised calibration sample, and on the C2 donor
realisation. They do not reflect training variability or calibration-sample
variability. We quantified the latter by resampling calibration patients (200
replicates), re-applying the frozen selection rule and recomputing the
estimates (\cref{sec:res-uncertainty}). Training variability is addressed only
by the replicate seed.

The registry specifies Holm correction across pathologies, and the original
SAP defines multiplicity families per hypothesis across the five
pathologies. The previous draft stated that per-pathology secondary analyses
used Holm correction, but no per-pathology analysis had been run. The Stage 10
code applied Holm to the four H4 bootstrap $p$-values, all of which were at the
bootstrap floor. The aggregate confirmatory contrasts (H1, H2-C1 and H3 for two
models) had no multiplicity adjustment. As a post hoc sensitivity we report
Bonferroni-adjusted intervals (level $1-0.05/6$) for those six contrasts, and
Holm-adjusted $p$-values across pathologies for the pathology-specific
contrasts.

\subsection{Post hoc analyses (Revision R1, 6 October 2026)}
\Cref{tab:status} lists every analysis with its status. Post hoc analyses are
exploratory and none is described as confirmatory. They were chosen to
answer specific validity questions: what the endpoint measures; whether the
label-source result reflects labels, cohort or thresholds; whether the
confidence score ranks errors; how selection and prediction changes combine
under interventions; whether the C2 constraints held; and how sensitive the
cross-site estimate is to defensible choices. No threshold, cutoff or score was
fitted on external data. Comparator selective scores were fitted on the source
calibration tier only.

\begin{table}[t]
\centering\small
\caption{Status of every reported analysis. ``Pre-results'' means specified and
implemented before the external point estimates were inspected on 7 August 2026.}
\label{tab:status}
\begin{tabular}{p{0.52\linewidth}p{0.4\linewidth}}
\toprule
Analysis & Status\\
\midrule
H1, H2 (C1), H4 contrasts; decision rule; bootstrap & registered 18 July; rule and bootstrap code pre-results (6 Aug); contrast code 7 Aug, after inspection\\
H3 as transfer-gap difference; H2 under C2 & defined in code after inspection (post hoc operationalisation; descriptive)\\
Coverage 70\%/90\%; $T=5$, $T=10$; findings labels (S1) & registered sensitivities; S1 implementation choices (re-selected thresholds, restricted source cohort) made after inspection\\
Second training seed (M1, M4) & decided 8 Aug, after inspection (not prespecified)\\
Estimator correction (evaluable-study) & correction of a verified defect (R1)\\
Label variants, aggregations, per-pathology contrasts, AURC/AUGRC, failure-detection AUROC & R1 post hoc (AURC, AUGRC, failure AUROC and pathology-level reporting were registered metrics or SAP requirements that had not been computed)\\
Controlled label-source decomposition; context selection/prediction split; C2 audit; calibration-sample uncertainty; comparator policies; view and demographic strata; additional C2 permutations; M2/M3 replicate & R1 post hoc, exploratory\\
\bottomrule
\end{tabular}
\end{table}
